\documentclass[aspectratio=169,10pt]{beamer}
\usepackage[utf8]{inputenc}
\usepackage[T1]{fontenc}
\usepackage[brazilian]{babel}
\usepackage{amsmath,amssymb,amsthm}
\usepackage{booktabs}

%\usetheme{Madrid}
\usecolortheme{default}
\setbeamertemplate{navigation symbols}{}
\setbeamertemplate{theorems}[numbered]

% Macros
\newcommand{\conj}[1]{\{#1\}}
\newcommand{\teto}[1]{\lceil #1\rceil}
\newcommand{\Abs}[1]{\left|#1\right|}
\newcommand{\E}[1]{\mathbb{E}\!\left[#1\right]}
\newcommand{\Var}[1]{\mathrm{Var}\!\left[#1\right]}
\newcommand{\ProbE}[1]{\Pr\!\left[#1\right]}
\newcommand{\paren}[1]{\left(#1\right)}
\newcommand{\coloneqq}{\mathrel{:=}}

\title[Streaming e o contador de Morris]{Algoritmos de \textit{streaming}\\e o contador de Morris}
\subtitle{Algoritmos Probabilísticos}
\author{J.D.}
\date{UFABC}

\begin{document}

\begin{frame}
  \titlepage
\end{frame}

\begin{frame}{Fluxos de dados}
  Um \alert{fluxo} é uma sequência extremamente longa de itens de um
  universo, lida \alert{uma única vez}, em ordem. Há dados demais para
  armazenar e eles chegam rápido demais.
  \medskip
  \begin{itemize}
  \item \textbf{Roteador de rede:} quantos endereços de origem distintos
    aparecem? Quais enviam mais tráfego?
  \item \textbf{Servidor web / busca:} quantos visitantes únicos? Quais
    os termos mais buscados?
  \item \textbf{Sensores:} leituras contínuas, em dispositivos com
    memória muito limitada.
  \end{itemize}
\end{frame}

\begin{frame}{O modelo de \textit{streaming}}
  Os dados chegam um a um:
  \[
    \sigma=(x_1,x_2,\dots,x_m),\qquad x_i\in\conj{1,\dots,n},
  \]
  com $m$ muito grande e possivelmente desconhecido.
  \begin{itemize}
  \item Memória muito menor que $m$, idealmente
    \alert{polilogarítmica}: $(\log m+\log n)^{O(1)}$ bits.
  \item Resultados exatos costumam ser impossíveis. Recorremos a
    \alert{aleatoriedade} e a aproximações $(\varepsilon,\delta)$:
    com probabilidade $\geq1-\delta$, o erro relativo é $\leq\varepsilon$.
  \end{itemize}

  \medskip

  Para comparação, memorizar o fluxo inteiro custaria
  $m\log n$ bits, e manter um vetor com um bit para cada elemento do
  universo, marcando os que já apareceram, custaria $n$ bits.
\end{frame}

\begin{frame}{Problemas básicos}
  \begin{enumerate}
  \item[a)] \textbf{Elementos distintos:} quantos valores diferentes
    aparecem?
  \item[b)] \textbf{Estimativa de frequências:} quantas vezes $i$
    apareceu, sem um contador por elemento (\textit{heavy hitters}).
  \item[c)] \textbf{Momentos de frequência:}
    $F_k=\sum_i f_i^k$. Aqui $F_0$ = distintos, $F_1=m$, $F_2$ mede a
    não uniformidade.
  \item[d)] \textbf{Amostragem em fluxo:} amostra uniforme de $\ell$
    elementos de um fluxo de tamanho desconhecido.
  \item[e)] \textbf{Contagem aproximada:} contar $m$ itens com muito
    menos que $\log_2 m$ bits.
  \end{enumerate}
  \medskip
  Hoje: o item e), o mais simples, devido a Morris (1978).
\end{frame}

\begin{frame}{Contagem exata}
  Fluxo de comprimento em $\conj{1,\dots,m}$, com $m$ conhecido.
  \begin{itemize}
  \item Mantemos um contador $c$, incrementado a cada item.
  \item Basta guardar $c \bmod m$: são $m$ valores, ou seja,
    $\teto{\log_2 m}$ bits. O resto $0$ é lido como $m$.
  \end{itemize}
  \medskip
  \alert{É possível usar bem menos?} Para algoritmos determinísticos
  exatos, não.
\end{frame}

\begin{frame}{Limite inferior para contagem exata}
  \begin{block}{Proposição}
    Seja $m\geq2$. Nenhum algoritmo determinístico com no máximo
    $\teto{\log_2m}-1$ bits devolve o comprimento exato de toda
    sequência de comprimento em $\conj{1,\dots,m}$. O mesmo vale para
    algoritmos aleatorizados que nunca erram.
  \end{block}
  \textbf{Ideia da prova.}
  \begin{itemize}
  \item Com $b=\teto{\log_2m}-1$ bits há $2^b<m$ estados.
  \item Seja $S_i$ o estado após ler $(1,\dots,1)$ de comprimento $i$,
    $i=1,\dots,m$. Pelo princípio da casa dos pombos, $S_i=S_j$ para
    algum $i\neq j$.
  \item A saída é função do estado final: mesma resposta, mas as
    corretas ($i$ e $j$) diferem.
  \item Aleatorizado que nunca erra: fixada a semente, torna-se
    determinístico.
  \end{itemize}
  O limite é justo: o contador módulo $m$ usa $\teto{\log_2m}$ bits.
\end{frame}

\begin{frame}{$(\varepsilon,\delta)$-aproximação}
  \begin{block}{Definição}
    Dados $\varepsilon,\delta\in(0,1)$, um algoritmo aleatorizado é uma
    \alert{$(\varepsilon,\delta)$-aproximação} se, para todo fluxo de
    comprimento $m$, calcula $\widehat m$ tal que
    \[
      \ProbE{\Abs{\widehat m-m}\leq\varepsilon m}\geq1-\delta,
    \]
    com a probabilidade sobre as escolhas aleatórias do algoritmo.
  \end{block}
  \begin{itemize}
  \item $\varepsilon$ é a \alert{precisão}; $1-\delta$ é a
    \alert{confiança}.
  \item Quanto menores $\varepsilon$ e $\delta$, melhor a garantia e
    maior a memória.
  \end{itemize}
\end{frame}

\begin{frame}{O algoritmo de Morris}
  \begin{block}{Algoritmo}
    \textbf{Entrada:} fluxo $(x_1,\dots,x_m)$ (os valores $x_i$ são
    ignorados). \\ \textbf{Saída:} estimativa $\widehat m$ de $m$.\\[2pt]
    $X\gets0$\\
    \textbf{para cada} item lido: \\
    \phantom{\textbf{para cada}} com probabilidade $2^{-X}$, faça
    $X\gets X+1$\\
    \textbf{devolva} $2^X-1$
  \end{block}\pause
  \textbf{Intuição.}
  \begin{itemize}
  \item Com $X=j$, o número de itens até o incremento é geométrico, de
    média $2^j$.
  \item Chegar a $X=j$ leva em média $1+2+\dots+2^{j-1}=2^j-1$ itens:
    daí a saída $2^X-1$.
  \item Só guardamos $X\approx\log_2m$, ou seja, $O(\log\log m)$ bits
    (com alta probabilidade).
  \end{itemize}
\end{frame}

\begin{frame}{Recorrência do contador}
  Seja $X_i$ o valor de $X$ após $i$ itens, $X_0=0$. Dado $X_{i-1}=j$:
  \[
    X_i=
    \begin{cases}
      j+1 & \text{com probabilidade } 2^{-j},\\
      j & \text{com probabilidade } 1-2^{-j}.
    \end{cases}
  \]
  Para $f\colon\mathbb N\to\mathbb N$, defina
  $g_f(j)\coloneqq(1-2^{-j})f(j)+2^{-j}f(j+1)$. Pela lei da esperança
  total,
  \[
    \E{f(X_i)}=\sum_j\Pr[X_{i-1}=j]\,g_f(j)=\E{g_f(X_{i-1})}.
    \tag{$\ast$}
  \]
  A soma é finita, pois $X_i\in\conj{0,\dots,i}$.
\end{frame}

\begin{frame}{O estimador é não viesado}
  \begin{block}{Teorema 1}
    \textit{ Para toda entrada de comprimento $m\geq0$,
      $\E{2^{X_m}}=m+1$. Logo $\widehat m=2^{X_m}-1$ satisfaz
      $\E{\widehat m}=m$.}
  \end{block}

  \pause
  
  \textbf{Prova.} Indução em $m$. Para $f(j)=2^j$,
  \[
    g_f(j)=(1-2^{-j})2^j+2^{-j}2^{j+1}=2^j+1.
  \]
  \pause
    Por $(\ast)$ e pela hipótese de indução,
  \[
    \E{2^{X_m}}=\E{2^{X_{m-1}}}+1=m+1.\qquad\square
  \]
\end{frame}

\begin{frame}{Segundo momento e variância}
  
  \begin{block}{Lema 1}
    $\E{(2^{X_m})^2} = \E{4^{X_m}} = \tfrac32m^2+\tfrac32m+1$.
  \end{block}

  \pause

  
  \textbf{Prova.} Para $f(j)=4^j$: $g_f(j)=4^j+3\cdot2^j$. Por $(\ast)$
  e pelo teorema anterior,
  \[
    \E{4^{X_i}}=\E{4^{X_{i-1}}}+3\E{2^{X_{i-1}}} =
    \E{4^{X_{i-1}}}+3i 
    \;\Rightarrow\;
    \E{4^{X_m}}=1+3\sum_{i=1}^m i=\tfrac32m^2+\tfrac32m+1.
  \]
  
  \pause

  \begin{block}{Corolário 1}
    $\Var{\widehat m}=\E{(2^{X_m})^2}-  ( \E {2^{X_m}})^2 =
    \dfrac{m(m-1)}2.$
  \end{block}
\end{frame}

\begin{frame}{Um único contador não basta}
  Para $m\geq1$, $\Var{\widehat m}\leq m^2/2$: o desvio padrão é da ordem
  de $m/\sqrt2$, comparável a $m$. Por Chebyshev,
  \[
    \ProbE{\Abs{\widehat m-m}\geq\varepsilon m}
    \leq\frac{\Var{\widehat m}}{\varepsilon^2m^2}\leq\frac1{2\varepsilon^2}.
  \]
  \begin{itemize}
  \item Só é não trivial para $\varepsilon>1/\sqrt2$.
  \item Com $\varepsilon\in(0,1)$, o limite é sempre $>1/2$: nunca dá
    $\delta<1/2$.
  \end{itemize}
  \medskip\pause
  \alert{Precisamos melhorar a concentração:} dois truques.
  \begin{enumerate}
  \item \textbf{Média} de cópias independentes: divide a variância.
  \item \textbf{Mediana} de várias médias: amplifica a confiança.
  \end{enumerate}
\end{frame}

\begin{frame}{O truque da média}
  Executamos $t$ cópias independentes do algoritmo sobre o mesmo fluxo e
  devolvemos
  \[
    \widehat m^*=\frac1t\sum_{i=1}^t\widehat m_i .
  \]
  \begin{itemize}
  \item $\E{\widehat m^*}=m$ e, pela independência,
    $\Var{\widehat m^*}=\dfrac{\Var{\widehat m_1}}t\leq\dfrac{m^2}{2t}$.
  \item Chebyshev:
    $\ProbE{\Abs{\widehat m^*-m}\geq\varepsilon m}\leq\dfrac1{2t\varepsilon^2}$.
  \item Com $t=\teto{1/(2\varepsilon^2\delta)}$, obtém-se uma
    $(\varepsilon,\delta)$-aproximação.
  \end{itemize}
  \alert{Problema:} o custo cresce com $1/\delta$.
\end{frame}

\begin{frame}{O truque da mediana}
  \begin{itemize}
  \item $t=\teto{3/(2\varepsilon^2)}$: cada média erra com probabilidade
    $\leq1/3$ (Chebyshev).
  \item Repetimos $s$ vezes (com $s$ ímpar) e devolvemos a
    \alert{mediana} $\widehat m^{**}$ das $s$ médias.
  \item $Y_i$: indicador de que a $i$-ésima média é boa. Independentes,
    $\Pr[Y_i=1]\geq2/3$, logo $\E Y\geq2s/3$ para $Y=\sum Y_i$.
  \item Mediana ruim $\Rightarrow$ pelo menos $(s+1)/2$ médias ruins
    $\Rightarrow$ $Y-\E Y<-s/6$.
  \item Chernoff--Hoeffding:
    \[
      \ProbE{\Abs{\widehat m^{**}-m}\geq\varepsilon m}
      \leq\exp\paren{-\frac{2(s/6)^2}{s}}=e^{-s/18}.
    \]
  \end{itemize}
  Basta $s\geq18\ln(2/\delta)$ para erro $\leq\delta/2$.
\end{frame}

\begin{frame}{Por que Chebyshev na média e Chernoff na mediana?}
  \begin{itemize}
  \item As saídas $\widehat m_i=2^{X_m}-1$ são \alert{ilimitadas}; só
    controlamos bem a variância $\Rightarrow$ Chebyshev.
  \item Os indicadores $Y_i\in\{0,1\}$ são limitados e independentes
    $\Rightarrow$ Chernoff--Hoeffding.
  \item Só com Chebyshev na mediana, precisaríamos de $s=\Omega(1/\delta)$.
    Com Chernoff, $s=O(\log(1/\delta))$.
  \end{itemize}
\end{frame}

 
\begin{frame}{Memória fixa: reservar e abortar}
  \begin{itemize}
  \item O espaço usado é \alert{aleatório}: $X$ pode, raramente, passar
    do limiar.
  \item O algoritmo \alert{reserva} o espaço do limiar e \alert{aborta}
    se algum contador o exceder (probabilidade $\leq\eta$).
  \item Com $\eta=\delta/2$: metade de $\delta$ para o erro de
    precisão, metade para o aborto.
  \end{itemize}
  \medskip
  \textbf{Exemplo.} $m=10^9$, $N=1$, $\eta=10^{-3}$: o limiar é
  $10^{12}$ e $\log_2 10^{12}\approx39{,}9$. Com probabilidade
  $\geq99{,}9\%$, $X<40$, que cabe em \alert{6 bits}; o contador exato
  precisaria de $30$ bits.
\end{frame}

\begin{frame}{Garantias finais}
  Para $\varepsilon,\delta\in(0,1)$ com $1/(\varepsilon\delta)\leq m^{O(1)}$,
  existem $(\varepsilon,\delta)$-aproximações de $m$ com memória:
  \medskip
  \begin{center}
    \begin{tabular}{lll}
      \toprule
      Variante & Contadores & Bits \\
      \midrule
      Média & $\teto{1/(\varepsilon^2\delta)}$ &
        $O\paren{\dfrac{\log\log m}{\varepsilon^2\delta}}$ \\[8pt]
      Mediana de médias & $ts$, $t=\teto{\tfrac3{2\varepsilon^2}}$,
        $s\geq18\ln\tfrac2\delta$ &
        $O\paren{\dfrac{\ln(2/\delta)\,\log\log m}{\varepsilon^2}}$ \\
      \bottomrule
    \end{tabular}
  \end{center}
  \medskip
  A mediana troca $1/\delta$ por $\ln(2/\delta)$: preferível quando
  $\delta$ é pequeno.
\end{frame}

\begin{frame}{Resumo}
  \begin{itemize}
  \item Fluxos: leitura única, memória polilogarítmica, aproximação
    aleatorizada $(\varepsilon,\delta)$.
  \item Contagem exata exige $\teto{\log_2 m}$ bits (limite inferior
    justo).
  \item Contador de Morris: $\E{2^{X_m}-1}=m$ e
    $\Var{2^{X_m}-1}=m(m-1)/2$, com $O(\log\log m)$ bits.
  \item \textbf{Média} reduz a variância; \textbf{mediana} amplifica a
    confiança.
  \end{itemize}
  \medskip
  \small R.~Morris, \textit{Counting large numbers of events in small
    registers}, Communications of the ACM 21(10), 840--842, 1978.
\end{frame}

\end{document}

%%% Local Variables:
%%% mode: LaTeX
%%% TeX-master: t
%%% End:
